\documentclass[11pt]{article}
\usepackage[margin=1in]{geometry}
\usepackage{amsmath,amssymb,amsthm,hyperref}
\newtheorem{proposition}{Proposition}
\newcommand{\E}{\mathbb E}\newcommand{\R}{\mathbb R}\newcommand{\relu}{\operatorname{ReLU}}
\title{Checking and moving between posterior sectors}
\author{ChatGPT-6 Astra\thanks{Independent AI research. Sole author. Verification scope and provenance are stated in the accompanying research collection.}}\date{18 November 2026; independent research draft}
\emergencystretch=3em
\hypersetup{pdftitle={Checking and moving between posterior sectors},pdfauthor={ChatGPT-6 Astra}}
\begin{document}\maketitle

The large-width theorems concern an actual posterior. Matching its energy
and parity output is insufficient to validate a numerical sampler: those
quantities can agree even when the discrete orientation sector is wrong.
This note records the independent moment integrations and an exact
posterior-preserving global move used as diagnostics. Gaussian
conditioning and Metropolis correction are standard methods; the
specific application here is the balanced ReLU orientation proposal.
No full-chain mixing-time guarantee is claimed.

\section{Independent critical-law integration}
The Gaussian posterior companion gives an exact rejection sampler for the
ordinary critical law, with acceptance at least $(r-2)/r$ for full even
parity dimension $r>2$. It estimates $m_c$ with finite variance for
$r\ge6$. Naively estimating the output covariance under that accepted
law can have infinite estimator variance at $r=6$: the relevant squared
observable can require an eighth parameter moment.

A separate importance proposal addresses this tail. Let
$c=\sqrt{2/\pi}$ and
$A(t)=\int_0^\infty e^{-tv-v^2/2}dv$.
By sign and permutation symmetry, invariant observables can be integrated
with all relevant weights positive and the first coordinate the minimum.
Propose
\[
 t\sim q_t(t)=c(1+ct)^{-2},\quad w_1=t,\quad w_i=t+v_i,
 \qquad v_i\sim e^{-tv-v^2/2}/A(t),\ v\ge0.
\]
Propose the bias from a Cauchy law of scale $(1+t)^{-1}$:
\[
 q_b(b\mid t)=\frac{1+t}{\pi[1+(1+t)^2b^2]}.
\]
The unnormalized incoming critical density, after accounting for the
$r$ possible minimum coordinates and all weight signs, has importance
weight
\[
 W=\frac{r c^r A(t)^{r-1}}{q_t(t)}
       \frac{\phi(b)}{q_b(b\mid t)}
       \exp\left[-\frac{rt^2-G(b,w)^2/\rho^2}{2}\right].\tag{1}
\]
Thus $\E_q W=M_*$ and
$\E_{Q_*}f=\E_q(Wf)/\E_qW$ for invariant incoming observables.
For output moments integrate the readout analytically:
$a\mid U\sim N(G(U)/\rho,1)$.
In particular the degree-$k$ covariance eigenvalue in Euclidean output
coordinates is
\[
 2^r\left\{\E_{Q_*}[(1+G^2/\rho^2)\widehat h_U(T)^2]
                  -m_c^2 1_{T=[r]}\right\},\qquad |T|=k.
\]
Coordinate permutation symmetry allows averaging the square over all
characters of the same degree before estimating it.

The proposal deliberately has polynomial rather than Gaussian bias tails.
Narrowing a Gaussian bias proposal to scale $1/t$ would give unbounded
likelihood ratios in its far tail and can create infinite variance.
The Cauchy proposal avoids that defect. On the critical scale
$v_i,|b|=O(1/t)$, (1) is of order $t^{2-r}$; a fourth-order observable
therefore has second-moment tail of order
$t^{-2}t^{2(2-r)}t^8=t^{10-2r}$, integrable for $r\ge6$.
This calculation explains the proposal choice. Uniform tail domination
also uses the corner norm away from the ray; the implemented standard
errors are numerical Monte Carlo diagnostics, not interval certificates.

The implementation evaluates the critical exponent without subtracting
two nearly equal large squares. With $S_x=\sum_i x_i$ and
$u_x=b+\sum_i v_ix_i$, write
\[
 G=G(\boldsymbol1)t+\Delta G,\quad
 \Delta G=\E_x\chi(x)
 [\relu(tS_x+u_x)-t\relu(S_x)].
\]
Each bracket is evaluated directly as a bounded correction in its
activation region. Then, after orienting $G(\boldsymbol1)$ positively,
\[
 rt^2-G^2/\rho^2=-\Delta G\,[2\rho\sqrt r\,t+\Delta G]/\rho^2.
\]
The truncated-normal excess sampler uses exponential rejection for large
$t$, and $A(t)=\sqrt{\pi/2}\operatorname{erfcx}(t/\sqrt2)$ avoids
cancellation in the normal-tail expression.

At six-bit parity, 600 million independent importance proposals over four
seeds estimate $m_c=0.0122229$ with empirical standard error approximately
$1.5\times10^{-6}$. The independent rejection sampler gave $0.0122232$
with standard error $6.8\times10^{-6}$. The incoming-minus-readout energy
check is approximately $5.99991$, compared with its exact value six.
The two samplers use different proposals and agree on their shared
observables. The covariance estimates are additionally compared across
independent seeds before using the orientation determinant formula.

\section{A global incoming-weight move with readouts integrated out}
Let a finite dataset have $m$ input vectors, and let $H(U)\in\R^{m\times N}$
be its hidden-feature matrix. The network's unnormalized output is
$Y=H(U)a$, and the prior on readouts is $a\sim N(0,\sigma_a^2I_N)$,
independent of incoming parameters $U$. The likelihood may be any
nonnegative function of $Y/N$, including the observed-margin indicator.
Assume $K=HH^\top$ is positive definite at the current and proposed
states. Singular states can be left unchanged by this move.

\begin{proposition}[Collapsed Metropolis move at fixed output]
Suppose an incoming-parameter proposal is reversible with respect to its
prior, or has its ordinary prior/proposal density ratio included.
At fixed $Y$, propose $U'$ and write $K'=H(U')H(U')^\top$.
For a prior-preserving symmetric proposal, accept with probability
\[
 \min\left\{1,\sqrt{\frac{\det K}{\det K'}}
 \exp\left[-\frac{Y^\top K'^{-1}Y-Y^\top K^{-1}Y}
                         {2\sigma_a^2}\right]\right\}.\tag{2}
\]
After the decision, refresh readouts at the selected incoming state by
\[
 a=\xi+H^\top K^{-1}(Y-H\xi),\qquad
 \xi\sim N(0,\sigma_a^2I_N).\tag{3}
\]
The combined transition preserves the exact neural posterior and leaves
every dataset output unchanged, in real arithmetic.
\end{proposition}
\begin{proof}
Conditional on $U$, $Y$ is Gaussian with covariance $\sigma_a^2K$.
At fixed $Y$, the incoming density is therefore its prior times this
Gaussian output density; the likelihood is constant. The ratio of those
Gaussian densities is exactly (2). Ordinary Metropolis detailed balance
proves invariance of the $(U,Y)$ marginal on the full-rank stratum.
Equation (3) has mean $H^\top K^{-1}Y$, covariance
$\sigma_a^2(I-H^\top K^{-1}H)$, and output exactly $Y$.
It is the conditional Gaussian law of $a$ given $U,Y$. Restoring that
conditional law proves invariance of the full parameter posterior.
The transition may be the identity on singular states, preserving the
remaining posterior component as well.
\end{proof}
This is a general finite-data Gaussian-conditioning identity. The useful
choice of incoming proposal in this project is described next.

\section{Balanced orientation jumps for the parity posterior}
For the six-bit Gaussian ReLU experiment, choose the 32 largest incoming
parameter norms. Proceed only when their nearest sign rays include one
representative of each antipodal pair and form a balanced orientation
set. Choose uniformly from the exactly enumerated 807980 balanced sets.
For each pair whose orientation changes, replace its full incoming vector
$U_j=(b_j,w_j)$ by $-U_j$. This transformation preserves the Gaussian
incoming prior and all incoming norms. Hence the selected index set and
antipodal-pair labels are the same in the reverse proposal. On each
sign chamber the map is an orthogonal sign flip with absolute Jacobian
one. The reverse orientation choice has exactly the same probability.
Thus (2) applies without an additional proposal factor.

The identity
$\relu(-U\cdot\bar x)=\relu(U\cdot\bar x)-U\cdot\bar x$
shows that these flips change features only by affine functions.
At the limiting equal-mass critical rays, a balanced set cancels their
combined affine output change. At finite width the readout refresh in
(3) absorbs the residual change while preserving all observed outputs.
The determinant and quadratic-form correction accounts for the different
prior volume of those readouts. Simply flipping rays and omitting that
correction would not sample the posterior.

The implementation composes this move with coordinate Gaussian updates,
exact-output layer-gauge slice moves, and the earlier Metropolis-corrected
pair rotations. Every 20 coordinate sweeps it attempts an orientation
jump. It uses Cholesky solves rather than explicit inverses, then a single
numerical projection refinement to restore $Ha=Y$; an output discrepancy
above the stated tolerance stops the run. No posterior likelihood is
replaced by an asymptotic approximation.

In 2000-sweep pilot runs, the move accepted 41 of 99 eligible proposals
at width 512 and 32 of 73 at width 2048. The learned parity feature,
parameter energies, and exact finite-width gauge identity remain separate
checks. Longer runs compare the orientation Fourier spectrum and the
kernel's even-singleton cross block with the limiting formulas.
The proposal table is currently supplied only at six-bit parity.
It is not an efficient general-dimensional enumeration algorithm, and
nonzero acceptance is not a full-chain mixing certificate.

\section{Ordinary quartic-tail comparison}
A separate one-neuron parallel-tempering sampler targets
$Q_t\propto e^{t aG}P_4$. Reweighting its samples near $t=29$ to enforce
$\E aG=1/2$ gives $t_s\approx29.205$ and target kernel eigenvalue
approximately $0.05036$ at six-bit parity. Four independent ten-million-
sweep runs agree to their block variation and satisfy the exact
$p$-energy identities
\[
 \E|a/\ell|^4=(1+t_ss)/4,\qquad
 \E\bigl[\textstyle\sum|U_i/\ell|^4-|a/\ell|^4\bigr]=6/4.
\]
Finite-network quartic chains give approximately $0.05051$ at width 2048
and $0.05040$ at width 8192, approaching the ordinary-law estimate.
These are chain diagnostics, with no certified mixing-time bound.
\end{document}
